\chapter{対数関数と指数関数}

この章では指数関数と対数関数を定義していきます。
指数関数は、高校数学ではとても当たり前に使用していましたが、$a^x$という「実数の実数乗」というものが一体なんなのかよくわかりません。
説明としては、既に実数の有理数乗は定義しているので、無理数$x$に収束する任意の有理数列$\{q_n\}_{n=1}^{+\infty}$をとってきて、
\[
  a^x:=\lim_{n\to+\infty}a^{q_n}
\]
で定義しようというものです。

これが間違いというわけではないですが、いきなり「これで定義しました」というのは、僕個人としては納得感が薄いというものです。
なんでかというと、これで定義しましたとなっても、
\begin{itemize}
  \item 右辺の極限は、$x$に収束する「任意の有理数列」$\{q_n\}$に対して、一意の実数に収束するのか？
  \item 指数法則$a^{x+y}=a^xa^y$を満たすか？
  \item 連続な逆関数をもつための条件(単調増加・単調減少)を満たすか？
  \item 微分は$a^x\log(a)$になるか？
  \item 逆関数$\log_a(x)$は微分可能か？
\end{itemize}
といったことを理路整然と説明するのが難しそうだからです。

そこで僕たちは、あえての対数関数から分け入ります。
\[
  \log(x)=\int_1^x \frac1\xi d\xi
\]
を知っている人には都合がいいことに、まだ$1/x$の積分が定義できていません。
これは$(0,+\infty)$で明らかに連続なので、積分可能であるはずです。
その積分値を返す関数として、対数関数を定義するのはすごく早いです。
早いというのは、対数関数は積分で定義しているので、対数関数が微分可能なことがすぐに出てきます。
そのほかの対数関数の満たすべき性質
\[
  \log(xy)=\log(x)+\log(y)
\]
なども比較的簡単に出てきます。
あとは中身を見ていきましょう。



\section{対数関数と指数関数}

$n\neq-1$のとき
\[
  \int x^ndx=\frac{x^{n+1}}{n+1}+C
\]
なのでした。
これは$n=-1$で成り立たない公式です。しかし
\[
  \frac1x:(0,+\infty)\to\mathbb{R}
\]
は連続なので、$(0,+\infty)$に含まれる任意の閉区間で積分可能なはずです。
これがどんな性質をもつ関数になるか調べていきましょう。

\subsection{対数関数}

$x>0$に対して
\[
  f(x):=\int_1^x\frac{1}{\xi}d\xi
\]
とおきます\footnote{
  積分範囲を$0$から取らず$1$から取っているのは、$x=0$だと$1/x$が定義されないからです。
  じゃあなんで$1$を取ったのかというと、そこはまあなんでもいいんですけど、ようは$f(x)$は対数関数なので$f(1)=0$としたいからです。
}。
この関数の性質を調べていきましょう。
まず微分すると
\[
  f'(x)=\frac1x>0
\]
であることから、$f$は単調増加です。
$f:(0,+\infty)\to\mathbb{R}$は特に単射です。
実は全射でもあることが、次の補題から理解できます。
\begin{lemma}
  任意の有理数$a$と実数$x>0$に対して、
  \[
    f(x^a)=af(x)
  \]
  が成り立つ。
\end{lemma}
\begin{proof}
  合成関数の微分から
  \[
    (f(x^a))'=(x^a)'f'(x^a)=ax^{a-1}\frac1{x^a}=\frac{a}{x}
  \]
  一方、微分の線形性から
  \[
    (af(x))'=af'(x)=\frac{a}{x}
  \]
  となり、
  \[
    (f(x^a))'=(af(x))'\iff (f(x^a)-af(x))'=0
  \]
  が成り立つ。
  微分して0になる関数は定数であるから、
  \[
    f(x^a)-af(x)=C
  \]
  とおくと、$x=1$のとき
  \begin{align*}
    f(1^a)&=f(1)=\int_1^{1}\frac1\xi d\xi=0\\
    af(1)&=af(1)=a\cdot\int_1^{1}\frac1\xi d\xi=a\cdot0=0
  \end{align*}
  ゆえに$C=f(1^a)-af(1)=0$となり、題意が従う。
\end{proof}
一方$1\leq\xi\leq2$のとき$1/2\leq\xi\leq1$なので、
\[
  f(2)=\int_1^2 \frac1\xi d\xi \geq \int_1^2 \frac12 d\xi=\frac12>0
\]
を満たします。
そこで、任意の$y>0$に対して、$nf(2)>y$となる自然数$n$を取ることができ、補題から
\[
  nf(2)=f(2^n)>y>0=f(1)
\]
ゆえに中間値の定理より、ある$1<c<2^n$が存在して、$f(c)=y$を満たします。

任意の負の実数になることも、$1/2\leq\xi\leq1$のとき$1\leq\xi^{-1}\leq2$なので
\[
  f\left(\frac12\right)=\int_1^{1/2}\frac1\xi d\xi=-\int_{1/2}^1\frac1\xi d\xi\leq-\int_{1/2}^11d\xi=-\frac12<0
\]
となることから、正の実数と同様に示すことができます。

またこの関数$f(x)=\int_1^x\xi^{-1}d\xi$は、次のような意味で「掛け算を足し算に変換」します。
\begin{theorem}
  任意の実数$x>0$, $y>0$に対して、
  \[
    f(xy)=f(x)+f(y)
  \]
  が成り立つ。
\end{theorem}
\begin{proof}
  両辺を$y$を定数とみて$x$で微分すると、合成関数の微分から
  \begin{align*}
    \frac{df(xy)}{dx}=\frac{d}{dx}\int_1^{xy}\frac1\xi d\xi=y\cdot\frac1{xy}=\frac1x\\
    \frac{d}{dx}(f(x)+f(y))=\frac{df(x)}{dx}=\frac1x
  \end{align*}
  となり、積分定数を除いて$f(xy)=f(x)+f(y)$が成り立ちます。
  一方で$x=1$のときに両辺を比較すれば、厳密に等式がなりたつことがわかる。
\end{proof}

このように、存在だけが分かっている関数も、だんだん素性がわかってきました。
といったところで、マシな名前をつけてあげましょう。
\begin{definition}[対数関数]
  \[
    \log:(0,+\infty)\to\mathbb{R};x\mapsto\int_1^x\frac1\xi d\xi
  \]
  を対数関数と呼ぶ。
\end{definition}

対数関数に関する重要な性質をまとめると次になります。
\begin{itemize}
  \item $(0,+\infty)$から$\mathbb{R}$への単調増加かつ全射な関数
  \item $(\log(x))'=1/x$ (${}^\forall x>0$)
  \item $\log(x^a)=a\log(x)$ (${}^\forall x>0$, ${}^\forall a\in\mathbb{Q}$)
  \item $\log(1)=0$
  \item $\log(xy)=\log(x)+\log(y)$ (${}^\forall x,y>0$)
  \item $\log(1/x)=-\log(x)$ ($\log(x^a)=a\log(x)$において$a=-1$とする)
\end{itemize}
後半3つは、$\log$が、「区間$(0,+\infty)$と積がなす群」から「$\mathbb{R}$と和がなす群」への群準同型になっていることを意味します。
微積分とはあまり関係のなさそうな性質ではありますが、ものごとはいろんな見方ができた方が、悪いってことはないです。

\subsection{指数関数}

対数関数は区間$(0,+\infty)$から$\mathbb{R}$への単調増加な全射なので、逆写像が存在します。
これを仮に
\[
  g:\mathbb{R}\to(0,+\infty)
\]
とおきます。つまり
\begin{align*}
  \log(g(x))&=x\quad({}^\forall x\in\mathbb{R})\\
  g(\log(x))&=x\quad({}^\forall x>0)
\end{align*}
が成り立つということです。

定義と対数関数の性質からすぐにわかるのは、
\begin{itemize}
  \item $g(0)=1$
  \item $g(x+y)=g(\log(g(x))+\log(g(y)))=g(\log(g(x)g(y)))=g(x)g(y)$
  \item $g(-x)=g(-\log(g(x)))=g(\log(1/g(x)))=1/g(x)$
\end{itemize}
あたりのことです。
つまり、$g$は「$\mathbb{R}$と和がなす群」から「区間$(0,+\infty)$と積がなす群」への群準同型になっているということです。

また、$a\in\mathbb{Q}$に対して$\log(x^a)=a\log(x)$が成り立っていたので、$x=g(1)$を代入すると
\begin{align*}
  \log(g(1)^a)=a\log(g(1))=a
\end{align*}
$g$にこの値を代入すると、
\[
  g(a)=g(\log(g(1)^a))=g(1)^a
\]
となります。
この等式は任意の$a\in\mathbb{Q}$に対して成り立つから、$g(x)$の$x\in\mathbb{Q}$での値は$g(1)$の値のみによって決まると言えます。
そこで
\begin{definition}
  実数$e:=g(1)$を\textbf{自然対数の底}または\textbf{ネイピア数}などと呼ぶ。
\end{definition}
\noindent
ことにして、任意の$x\in\mathbb{R}$に対して
\[
  e^x:=\exp(x):=g(x)
\]
と書きます\footnote{
  まだ実数の「無理数乗」を定義していないので、任意の$x\in\mathbb{R}$に対して$e\in\mathbb{R}$の$e^x$というのを見た目通りに解釈することは、まだできません。
  $e^x$は現状、既に存在が保証されている関数$g(x)$の別名にすぎないことに注意です。
}。

あと気になることは
\begin{itemize}
  \item $e^x$は連続か？
  \item $e^x$は微分可能か？できるとしたらどうなる？
\end{itemize}
あたりだと思います。
まずは連続であることを気合で示します。
\begin{theorem}
  $e^x$は連続関数である。
\end{theorem}
\begin{proof}
  任意の$a\in\mathbb{R}$に対して、
  \[
    \lim_{x\to a}e^x=e^a
  \]
  を示したい。任意の$\varepsilon>0$に対して、
  \[
    \delta:=\log(1+\varepsilon e^{-a})
  \]
  とおけば、
  \begin{align*}
    &|x-a|<\delta\\
    \iff&-\log(1+\varepsilon e^{-a})<x-a<\log(1+\varepsilon e^{-a})\\
    \iff&\log\frac1{1+\varepsilon e^{-a}}<x-a<\log(1+\varepsilon e^{-a})\\
    \iff&\frac1{1+\varepsilon e^{-a}}<e^{x-a}<1+\varepsilon e^{-a}\\
    \iff&\frac1{1+\varepsilon e^{-a}}-1<e^{x-a}-1<\varepsilon e^{-a}\\
    \iff&-\frac{\varepsilon e^{-a}}{1+\varepsilon e^{-a}}<e^{x-a}-1<\varepsilon e^{-a}\\
    \implies&-\varepsilon e^{-a}<e^{x-a}-1<\varepsilon e^{-a}\\
    \iff&|e^{x-a}-1|<\varepsilon e^{-a}
  \end{align*}
  となる\footnote{
    下から2つ目の$\implies$では、$1+\varepsilon e^{-a}>1$すなわち$1/(1+\varepsilon e^{-a})<1$ゆえに$-1/(1+\varepsilon e^{-a})>-1$であることを使っています。
  }。
  さらに両辺に$e^a$を掛ければ
  \[
    e^a|e^{x-a}-1|=|e^x-e^a|<\varepsilon
  \]
  が成り立つ。
  すなわち
  \[
    \lim_{x\to a}e^x=e^a
  \]
\end{proof}

連続性を使って、微分を求めましょう。
そのためには
\[
  \lim_{h\to0}\frac{e^{x+h}-e^x}{h}
\]
が存在するかどうかを調べる必要があるのですが、
\[
  \frac{e^{x+h}-e^x}{h}=e^x\frac{e^h-1}{h}
\]
なので、
\[
  \lim_{h\to0}\frac{e^h-1}{h}
\]
を求めれば分かります\footnote{
  これは$e^x$の$x=0$における微分係数に他ならないです。
}。

直感的には、まず$h=\log(1+k)$と変換すると、
\[
  \lim_{h\to0}\frac{e^h-1}{h}=\lim_{k\to0}\frac{k}{\log(1+k)}
\]
です。
一方で対数関数の定義から
\[
  (\log(x))'(x=1)=\frac{1}{1}=1
\]
なので、微分係数の定義に戻れば
\[
  1=\lim_{k\to0}\frac{\log(1+k)-\log(1)}{k}=\lim_{k\to0}\frac{\log(1+k)}{k}
\]
ゆえに、逆数の極限
\[
  \lim_{h\to0}\frac{e^h-1}{h}=\lim_{k\to0}\frac{k}{\log(1+k)}
\]
も1になるでしょう。
気合を入れて証明します。
\begin{theorem}
  \[
    \lim_{h\to0}\frac{e^h-1}{h}=1
  \]
  ゆえに、任意の$x\in\mathbb{R}$に対して
  \[
    (e^x)'=e^x
  \]
\end{theorem}
\begin{proof}
  $\varepsilon>0$を任意の正の実数とする。
  \[
    (\log(x))'(x=1)=1
  \]
  であるから、ある$\delta_1>0$が存在して、$0<|k|<\delta_1$のとき
  \[
    \left|\frac{\log(1+k)}{k}-1\right|<\min\left\{\frac12,\frac{\varepsilon}{2}\right\}
  \]
  を満たす。このとき、
  \begin{align*}
    &\left|\frac{\log(1+k)}{k}-1\right|<\min\left\{\frac12,\frac{\varepsilon}{2}\right\}\\
    \iff&1-\min\left\{\frac12,\frac{\varepsilon}{2}\right\}<\frac{\log(1+k)}{k}<1+\min\left\{\frac12,\frac{\varepsilon}{2}\right\}\\
    \implies&\frac12<\frac{\log(1+k)}{k}<\frac32\\
    \implies&\left|\frac{\log(1+k)}{k}\right|>\frac12
  \end{align*}
  であるから、
  \begin{align*}
    \left|\frac{k}{\log(1+k)}-1\right|=\frac{|1-\frac1k\log(1+k)|}{|\frac1k\log(1+k)|}<2\left|1-\frac{\log(1+k)}{k}\right|\leq\varepsilon
  \end{align*}
  また、この$\delta_1>0$に対して、$e^x$は連続関数であるから
  \[
    \lim_{h\to0}e^h=e^0=e^{\log(1)}=1
  \]
  ゆえ、ある$\delta_2>0$が存在して、$0<|h|<\delta_2$ならば
  \[
    |e^h-1|<\delta_1
  \]
  を満たす。
  \[
    k=e^h-1 \iff h=\log(1+k)
  \]
  とおくと、$0<|h|<\delta_2$ならば$0<|k|=|e^h-1|<\delta_1$であり\footnote{
    $k=0 \iff e^h=1 \iff \log(e^h)=\log(1) \iff h=0$に注意です。
  }、
  \begin{align*}
    \left|\frac{e^h-1}{h}-1\right|=\left|\frac{k}{\log(1+k)}-1\right|<\varepsilon
  \end{align*}
\end{proof}

指数関数の重要な性質をまとめると以下のようになります\footnote{
  $(e^x)^y=e^{xy}$も成り立ちますが、現在は$(e^x)^y$が任意の$y\in\mathbb{R}$に対して定義されていないので証明不可能です。
}。
\begin{itemize}
  \item $e^0=1$
  \item $e^{x+y}=e^xe^y$
  \item $e^{-x}=1/e^x$
  \item $(e^x)'=e^x$
\end{itemize}

特に指数関数には、微分しても姿を変えないという著しい性質があります。
これを読み替えれば、指数関数の原始関数も自分自身です(積分定数の違いを除いて)。
\[
  \int e^x dx = e^x + C
\]



\section{実数べき関数の定義とその微積分}

対数関数
\[
  \log(x):=\int_1^x \frac1\xi d\xi
\]
の逆関数として、指数関数
\[
  e^x
\]
が定義でき、しかも微分可能かつ微分可能であることが示せました。
これによって「実数の無理数乗」を自然に定義することができます。

\subsection{実数べき関数の定義と性質}

まず任意の$x\in\mathbb{R}$に対して、既に$e^x$が定義されています。
また
\begin{align*}
  e^{\log(x)}=x
\end{align*}
だったのと、まだ証明できていない
\[
  (e^x)^y=e^{xy}
\]
との整合性を考えれば、次のように定義するのが良いでしょう。
\begin{definition}
  $n\in\mathbb{R}\setminus\{0\}$に対して、関数
  \[
    (0,+\infty)\to\mathbb{R};x\mapsto e^{n\log(x)}
  \]
  を$x^n$と書く。
\end{definition}

こうであれば、保留していた次の主張を証明できます。
\begin{theorem}
  任意の$x,y\in\mathbb{R}$に対して、
  \[
    (e^x)^y=e^{xy}
  \]
\end{theorem}
\begin{proof}
  \[
    (e^x)^y=\exp(y\log(e^x))=\exp(yx)=e^{xy}
  \]
\end{proof}

また、この記法を用いるからには以下の指数法則が成り立ちます。
\begin{theorem}
  任意の$n,m\in\mathbb{R}$と$x>0$に対して、
  \[
    x^{n+m}=x^nx^m
  \]
\end{theorem}
\begin{proof}
  \[x^{n+m}=e^{(n+m)\log(x)}=e^{n\log(x)+m\log(x)}=e^{n\log(x)}e^{m\log(x)}=x^nx^m\]
\end{proof}

また、$x^n$の「$n$に関する連続性」も証明できます。
\begin{theorem}
  任意の$n\in\mathbb{R}\setminus\{0\}$と、$n$に収束する有理数列$\{\alpha_m\}_{m=1}^{+\infty}$に対して、
  \[
    x^n=\lim_{m\to+\infty}x^{\alpha_m}
  \]
  が成り立つ。
\end{theorem}
\begin{proof}
  $e^x$の連続性より、
  \[\lim_{m\to+\infty}x^{\alpha_m}=\lim_{m\to+\infty}e^{\alpha_m\log(x)}=e^{n\log(x)}=x^n\]
\end{proof}

これで$x^n=e^{n\log(x)}$という定義の妥当性も理解できたと思います。
微分や積分についても、計算問題にすぎませんが証明していきます。
\begin{theorem}
  任意の$n\in\mathbb{R}\setminus\{0\}$と$x>0$に対して、
  \[(x^n)'=nx^{n-1}\]
\end{theorem}
\begin{proof}
  合成関数の微分、および、微分の線形性より
  \[
    (x^n)'=(e^{n\log(x)})'=(n\log(x))'e^{n\log(x)}=\frac{n}{x}x^n=nx^{n-1}
  \]
\end{proof}

\begin{theorem}
  任意の$n\in\mathbb{R}\setminus\{0,-1\}$と$x>0$に対して、
  \[\int x^n dx=\frac{x^{n+1}}{n+1}+C\]
\end{theorem}
\begin{proof}
  右辺を微分すれば明らか。
\end{proof}

\subsection{$x^n$の定義域について}

上記の通り、任意の$n\in\mathbb{R}$に対して$x^n$は$x>0$で定義できます。
整数でない有理数$n$については、例えば$\sqrt{-1}$などは実数にならないため$x>0$で定義されるとみるのが普通です。
一方で、自然数$n>0$に対しては$x^n$は任意の$x\in\mathbb{R}$に対して定義できます。
負の整数$n<0$については、0除算を回避するため$x\neq0$で定義されます。

実は一番ややこしいのは$n=0$の時です。
$x\neq0$なら$x^0=1$なので、連続性の観点でみれば任意の$x\in\mathbb{R}$に対して、$x^0=1$となりそうです。
しかし、$x=0$のほうを固定して$n\neq0$が0に近づくことを考えると、
\[
  \lim_{n\to0}0^n=\lim_{n\to0}0=0
\]
となります。
つまり、連続性を保つべきという観点で見れば、
\[
  0^0
\]
がうまく定義されないのです。
\begin{example}
  \begin{align*}
    &\lim_{n\to+\infty}e^{-n}=0\\
    &\lim_{n\to+\infty}\frac1{n}=0\\
    &\lim_{n\to+\infty}\frac1{n^2}=0
  \end{align*}
  であるが、
  \begin{align*}
    &\lim_{n\to+\infty}(e^{-n})^{1/n}=\lim_{n\to+\infty}e^{-1}=\frac1e\\
    &\lim_{n\to+\infty}(e^{-n})^{1/n^2}=\lim_{n\to+\infty}e^{-1/n}=e^0=1
  \end{align*}
\end{example}
こうした事情もあり、どの文脈であっても$x^n$の定義域には気を付けるべきです。
常に定義域を明記するべきでしょう。

\subsection{一般の底の指数関数と対数関数}

上記のように$x$を固定して$n$を変数とみる関数$x^n$を考えることがあります。
見やすいように
\[
  a^x
\]
というものを考えるのです。
これを指数関数というのですが、
\begin{align*}
  a^x=(e^{\log(a)})^x=e^{x\log(a)}
\end{align*}
なので、これでもってそのまま定義としちゃいましょう。
\begin{definition}
  $a>0$に対して、
  \[
    a^x:\mathbb{R}\to\mathbb{R};x\mapsto e^{x\log(a)}
  \]
  を、\textbf{$a$を底とする指数関数}という。
\end{definition}

この定義によれば、$a>0$を底とする指数関数は微分可能な関数の合成関数ですから、次の主張が簡単に示せます。
\begin{theorem}
  任意の$a>0$、$x\in\mathbb{R}$に対して、
  \[
    (a^x)'=a^x\log(a)
  \]
\end{theorem}
\begin{proof}
  合成関数の微分より
  \[
    (a^x)'=(e^{x\log(a)})'=(x\log(a))'e^{x\log(a)}=\log(a)\cdot a^x=a^x\log(a)
  \]
\end{proof}

次に$a^x$の値域について調べましょう。
僕たちの定義によれば、最終的に$e^x$に合成するので、$a^x$の値域は正の実数です。
逆に$a\neq1\iff\log(a)\neq0$のとき、任意の正の実数$y$に対して、$x=\log(y)/\log(a)$とすれば
\[
  a^x=e^{x\log(a)}=e^{\log(y)}=y
\]
ゆえ、$a^x$は$\mathbb{R}$から$(0,+\infty)$への全射です。
さらに上記の微分から、$\log(a)$の符号によって単調増加・単調減少がかわります。
ただし$\log(a)=0 \iff a=1$のときは$a^x$は常に1を返す定値関数です。
以上をまとめて
\begin{theorem}
  任意の$a>0$と$x\in\mathbb{R}$に対して、$a^x>0$。
  また
  \begin{itemize}
    \item $a>1$のとき、$a^x$は単調増加
    \item $a=1$のとき、$a^x$は$1$への定値関数
    \item $0<a<1$のとき、$a^x$は単調減少
  \end{itemize}
  特に$a>0$かつ$a\neq1$のとき、$a^x$は$\mathbb{R}$から$(0,+\infty)$への全単射である。
\end{theorem}

しれっと流しましたが、上の議論で$a>0$かつ$a\neq1$のときに、$a^x$の逆関数を構成していました。
\begin{definition}
  $a>0$、$a\neq1$に対して、
  \[
    \log_a(x):(0,+\infty)\to\mathbb{R};x\mapsto\frac{\log(x)}{\log(a)}
  \]
  を、\textbf{$a$を底とする対数関数}という。
\end{definition}



\section{空気抵抗ありの場合の落下運動}

力学の話に戻りましょう。
指数・対数関数を利用すれば、\textbf{空気抵抗ありの場合の落下運動}を、運動方程式で解析することができます。

自由落下は、どちら向きに重力が働いているかによりますが、運動方程式は
\[
  m\ddot{\gamma}=(0,0,mg)
\]
で表されます。明らかに一次元での運動になるので、初めから$\gamma$は一次元の値をとるとし、
\[
  m\ddot{\gamma}=mg
\]
を考えましょう。

空気抵抗があるとき、それは$\gamma$の速さに比例する力が加わりそうという予感がします。
車の窓を開け、手を伸ばしてみると、車が速ければ速いほど、強い力を手に感じることができます(そんなことしちゃダメだぞ)。
従って、自由落下の運動方程式を次のように変えるのが良さそうです。
\[
  m\ddot{\gamma}=mg-k\dot{\gamma}
\]
$k>0$は比例定数です。
$k$は、定性的な結果があっていれば、観測事実から後から決める実数です。
いったん、速度がわかればよいという観点から
\[
  v(t):=\dot{\gamma}(t)
\]
とおくと、運動方程式は
\[
  \frac{dv}{dt}=g-\frac{k}{m}v
\]
となります。

このように何かの未知関数の微分が絡んだ方程式を、\textbf{微分方程式}といいます。
微分方程式の解法第一号は、今回使える\textbf{変数分離法}です。
$g-\frac{k}{m}v\neq0$の範囲で、
\begin{align*}
  \frac{dv}{dt}=g-\frac{k}{m}v &\iff \frac{1}{g-kv/m}\frac{dv}{dt}=1
\end{align*}
この両辺を$t$で積分するのですが、右辺は簡単に$t+C$になりますね。
左辺は、合成関数の積分を使って
\begin{align*}
  \int \frac{1}{g-kv/m}\frac{dv}{dt}dt &= \int \frac{1}{g-kv/m}dv=-\frac{m}{k}\log(g-kv/m)
\end{align*}
です。
積分定数がまだないですが、右辺にまだ不定積分が残ってるので、そちらに吸収させましょう。
ということで、まとめて
\[
  -\frac{m}{k}\log(g-kv/m)=t+C \iff  g-\frac{k}{m}v(t)=A\exp\left(-\frac{kt}{m}\right) \iff v(t)=\frac{mg}{k}-\frac{mA}{k}\exp\left(-\frac{kt}{m}\right)
\]
と求めることができます。
積分定数$A$は、例えば自由落下するボールを「静かに落とした」という初期条件$v(0)=0$を課すとすると、
\[
  0=\frac{mg}{k}-\frac{mA}{k} \quad \therefore \; A=g
\]
と求まるので、空気抵抗付きの自由落下ボールの運動は
\[
  v(t)=\frac{mg}{k}\left\{1-\exp\left(-\frac{kt}{m}\right)\right\}
\]
で記述できるというわけです。

$t$が十分大きいとき、$e^{-\alpha t}$はかなり急速に0に収束します。
自由落下ボールの速度も、十分な時間を掛ければ$v=mg/k$に急速に近づいていきます。
例えば雨粒が雲の高さから空気抵抗を無視して地表に落ちてくると、とんでもないスピードで落ちてきます。
空気抵抗のおかげで、人間のたんぱく質を貫かない程度の勢いにとどめてくれているんですね。
